\scanpage{54}
\mspage{1}
% archive foliation: 56
\uline{\unsure{11.1.84|16.1.84}}

\uline{Chirurgie conforme}

\textcircled{\unsure{1|A})} Surface conforme à bord. On \deleted{\unsure{peut donc}} le définit par des
cartes locales qui, \uline{au} voisinage d'un point du
bord, sont \unsure{modelées} sur \inserted{$\overline H=$} $(\Im z\geq0)\subset\mathbb C$.
Dans le cas $X$ orienté i.e. structure holomorphe, la
structure peut \deleted{se} définir par un faisceau
de $\mathbb C$-algèbres, qui sur le modèle
$\overline H$ précédent, aux points $x$ de $\delta\overline H=\mathbb R$, est \unsure{formé}
de germes de fonctions qui se \uline{prolongent} en
fonctions hol. \deleted{aux}\inserted{au} \uline{voisinage} de $x$ \uline{dans} $\mathbb C$.
À ne pas confondre avec le sous faisceau
plus grand du faisceau des fonctions
complexes continues, savoir \deleted{\unsure{celui des}}\inserted{les} fonctions
continues qui sont hol dans $\operatorname{Int}(X)$.

\uline{Théorème 1} La structure conforme de $X$ (surface
conforme à bord) est connue quand on
la connaît dans $\operatorname{Int}(X)$, (ce qui suffit à
déterminer le faisceau des fonctions \inserted{complexes} continues
qui sont conformes dans $\operatorname{Int}(X)$).

Il suffit de voir que cela permet de
reconstituer les \uline{cartes} conformes des voisinages
des points du bord. Je dis que ces cartes
